\clearpage
\scanpage{6}
\mspage{\unsure{\handwritten{(31}|(31}}

% Pencil cancellation across the typescript.
\newsavebox{\coteSixTypescript}
\begin{lrbox}{\coteSixTypescript}
\begin{minipage}[t]{.96\linewidth}
(les morphismes se définissant de la façon évidente), et on note qu'on a un homomorphisme évident
\[
\underset{(\mathrm Y,\handwritten{\xi})\handwritten{\unsure{\in|\mathrm E}}\mathrm C_{\handwritten{\varphi}}^{\handwritten{\unsure{o|0|\circ}}}}{\underset{\handwritten{\longrightarrow}}{\lim}}\ \overset{\handwritten{\scriptscriptstyle\vee}}{\mathrm Y}
\underset{\begin{tikzpicture}[x=1cm,y=1cm]\draw (.35,0)--(.38,.10);\end{tikzpicture}}{\mathrel{\mbox{-}\mkern-2mu\mbox{-}\mkern-2mu\mbox{-}\mkern-2mu\mbox{-}\mkern-2mu\mbox{-}\mkern-2mu\mbox{-}\mkern-3mu\handwritten{\to}}}\handwritten{\varphi}
\tag*{\llap{(\(\unsure{\underline{\overline{\ast}}|\underline{\ast}|\ast}\))\hspace{.83\linewidth}}}
\]
en remarquant que
\[
\operatorname{Hom}\bigl(\underset{\mathrm C_{\handwritten{\varphi}}}{\underset{\handwritten{\rightarrow}}{\lim}}\ \overset{\handwritten{\scriptscriptstyle\vee}}{\mathrm Y},\handwritten{\varphi}\bigr)
\handwritten{\simeq}\underset{\mathrm C_{\handwritten{\varphi}}}{\underset{\handwritten{\leftarrow}}{\lim}}\operatorname{Hom}\bigl(\overset{\handwritten{\scriptscriptstyle\vee}}{\mathrm Y},\handwritten{\varphi}\bigr)
\handwritten{\simeq}\underset{\mathrm C_{\handwritten{\varphi}}}{\underset{\handwritten{\leftarrow}}{\lim}}\handwritten{\varphi}(\overset{\handwritten{\sim}}{\mathrm Y}).
\]
Or on a un élément canonique dans $\underset{\mathrm C_{\handwritten{\varphi}}}{\underset{\handwritten{\leftarrow}}{\lim}}\handwritten{\varphi}(\overset{\handwritten{\sim}}{\mathrm Y})$, en associant à tout $(\mathrm Y,\handwritten{\xi})$ l'élément $\handwritten{\xi}$ de $\handwritten{\varphi}(\mathrm Y)$. Utilisant le fait que les $\underset{\handwritten{\leftarrow}}{\lim}$ finies existent dans $\mathrm C$, et que le foncteur $\mathrm Y\handwritten{\leadsto\varphi}(\overset{\handwritten{\sim}}{\mathrm Y})$ y commute, on voit que dans $\mathrm C_{\handwritten{\varphi}}$ les $\underset{\handwritten{\leftarrow}}{\lim}$ finies existent; à fortion $\mathrm C_{\handwritten{\varphi}}^{\handwritten{o}}$ est pseudo-filtrant. Utilisant que tout faisceau $\mathrm F$ est limite inductive des faisceaux de la forme $\overset{\handwritten{\sim}}{\mathrm Y}$, et utilisant le fait que $\handwritten{\varphi}$ commute aux limites inductives, on conclut aisement que l'homomorphisme de foncteurs ($\underline{\overline{\ast}}$) est bijectif, et donne donc la représentation annoncée.

\handwritten{7.9.2.} Remarquons également que si $\underline{\mathrm E}$ est un topos, $\mathrm Y$ un objet de $\underline{\mathrm E}$, d'où un topos $\underline{\mathrm E}_{/\mathrm Y}$, alors la donnée d'un foncteur fibre $\handwritten{\varphi}$ pour $\underline{\mathrm E}_{/\mathrm Y}$ équivaut à la donnée d'un couple $(\handwritten{\varphi},\handwritten{\xi})$, où $\handwritten{\varphi}$ est un foncteur fibre pour $\underline{\mathrm E}$, et $\handwritten{\xi\mathrel{\mathrm E}\varphi}(\mathrm Y)$. Au foncteur fibre $\handwritten{\psi}:\underline{\mathrm E}_{/\mathrm Y}\mathrel{\mbox{-}\mkern-2mu\mbox{-}\mkern-2mu\mbox{-}\mkern-3mu\handwritten{\to}}(\mathrm{Ens})$ on associe le couple $(\handwritten{\varphi},\handwritten{\xi})$, où $\handwritten{\varphi}$ est le composé $\handwritten{\varphi}(\mathrm Z)=\handwritten{\psi}(\mathrm Z\mathrm{x}\mathrm Y)$, $\handwritten{\xi\in\varphi}(\mathrm Y)=\handwritten{\psi}(\mathrm Y\mathrm{x}\mathrm Y)$ est l'image de l'unique élément de $\handwritten{\psi}(\mathrm Y)$ par le morphisme diagonal de $\mathrm Y\mathrm{x}\mathrm Y$.

\handwritten{7.9.3.} Revenons maintenant au cas du site $\mathrm C=\overset{\handwritten{\sim}}{\mathrm X_{\text{ét}}}$, et soit $\handwritten{\varphi}:\overset{\handwritten{\sim}}{\mathrm X_{\text{ét}}}\mathrel{\mbox{-}\mkern-2mu\mbox{-}\mkern-2mu\mbox{-}\mkern-3mu\handwritten{\to}}(\mathrm{Ens})$ un foncteur fibre sur le topos étale de $\mathrm X$.
\end{minipage}
\end{lrbox}
\noindent\begin{tikzpicture}
\node[inner sep=0pt,outer sep=0pt] (body) {\begingroup\renewcommand{\ULthickness}{0pt}\deleted{\usebox{\coteSixTypescript}}\endgroup};
\draw[gray,thin] ([xshift=.44\linewidth,yshift=.8em]body.south west) -- ([xshift=-.43\linewidth,yshift=-1.3em]body.north east);
\end{tikzpicture}
% archive sheet number: 5
